% Two legal binary search trees over the same seven keys. The arrival order is
% the only difference between them, and it already costs a factor of more than
% two in search depth at this size.
%
% Positions are explicit rather than delegated to TikZ's tree layout: the right
% tree has one child per node, and the automatic layout would stack those
% vertically instead of showing the list running down to the right.
\begin{tikzpicture}[x = 1mm, y = 1mm]

  % ------------------------------------------------ left: a balanced arrival
  \begin{scope}
    \node[tnode] (a4) at (  0,   0) {4};
    \node[tnode] (a2) at (-14, -13) {2};
    \node[tnode] (a6) at ( 14, -13) {6};
    \node[tnode] (a1) at (-21, -26) {1};
    \node[tnode] (a3) at ( -7, -26) {3};
    \node[tnode] (a5) at (  7, -26) {5};
    \node[tnode] (a7) at ( 21, -26) {7};
    \foreach \p/\c in {a4/a2, a4/a6, a2/a1, a2/a3, a6/a5, a6/a7}
      \draw[tedge] (\p) -- (\c);

    \node[tannot] at (0,  10) {inserted 4, 2, 6, 1, 3, 5, 7};
    \node[tannot] at (0, -36) {height 2, at most three comparisons};
  \end{scope}

  % -------------------------------------------- right: the keys arrive sorted
  \begin{scope}[xshift = 72mm]
    \foreach \k [count = \i from 0] in {1, 2, 3, 4, 5, 6, 7}
      \node[tnode] (b\k) at (\i * 6 - 18, -\i * 8) {\k};
    \foreach \p/\c in {b1/b2, b2/b3, b3/b4, b4/b5, b5/b6, b6/b7}
      \draw[tedge] (\p) -- (\c);

    \node[tannot] at (0,  10) {inserted 1, 2, 3, 4, 5, 6, 7};
    \node[tannot] at (0, -60) {height 6, up to seven comparisons};
  \end{scope}

\end{tikzpicture}
